\section{讲座背景与目标}
\subsection{讲者与课程定位}
这场讲座属于 UC Berkeley CS294-280 Sp25 进阶课程，由 Meta Research 的 Jason Weston 主持。讲座开场即指出当前 reasoning 研究正经历类似 01/R1 与 DeepSeek-R1 的跳跃：模型在多模态 benchmark 上表现迅猛，但每一次 leaps 都也暴露出思维错误无法被及时发现的问题。因此，本节旨在重构讲座的定位：它既是对 Chain-of-Thought 近年来发展史（2022 年才真正流行）的小结，也是在 ``让模型自己定义 reasoning 任务'' 与 ``让它自己判断结果'' 之间架桥。

\subsection{讲座旨趣}
Weston 把主题归结为 self-improving language models——模型在训练过程中逐步学习如何自己提问、自己评分、自己优化。讲者在开场多次提到 ``我觉得我们这一代 AI 研究者非常幸运，可以目睹这些飞速进步''，并强调 ``analogous to DeepSeek-R1 这样的 benchmark 事件，即便在短短几个月内也能看到 reasoning accuracy 的质变''。因此，笔记接下来会按照推理目标、训练闭环、评价体系和部署挑战四个维度逐步展开。

\begin{knowledgebox}{Chain-of-Thought 的时代背景}
Chain-of-Thought 训练框架从 2022 年开始爆发式流行，Weston 把它视为 ``reasoning 的快速起点''，但又提醒：真正的推理能力需要模型不仅输出推理链，还要能评估哪一条推理更靠谱。
\end{knowledgebox}

\subsection{本章小结}
本节通过历史节点（Chain-of-Thought 的崛起、DeepSeek-R1、01/R1）及主持人对 ``self-improvement'' 感想的阐述，确立了整场讲座的顶层目标和评价尺度。接下来需要关注的核心问题是：我们如何把 ``模型的自省能力'' 拆解成实际训练步骤。

